\documentclass[12pt]{article}
\usepackage{sbc-template}
\usepackage{graphicx,url}
\usepackage[utf8]{inputenc}
\usepackage[brazil]{babel}
\usepackage{placeins}

% Floats no estilo IEEE (coluna única): topo/rodapé, sem esticar páginas
\setcounter{topnumber}{2}
\setcounter{bottomnumber}{2}
\setcounter{totalnumber}{4}
\renewcommand{\topfraction}{0.9}
\renewcommand{\bottomfraction}{0.8}
\renewcommand{\textfraction}{0.08}
\renewcommand{\floatpagefraction}{0.85}
\setlength{\textfloatsep}{8pt plus 2pt minus 2pt}
\setlength{\floatsep}{8pt plus 2pt minus 2pt}
\setlength{\intextsep}{8pt plus 2pt minus 2pt}
\setlength{\abovecaptionskip}{6pt}
\setlength{\belowcaptionskip}{2pt}

\newcommand{\figwidth}{0.95\linewidth}
\newcommand{\includefig}[2][]{%
  \includegraphics[width=\figwidth,#1]{#2}}
\newcommand{\includefigtall}[1]{%
  \includefig[height=0.32\textheight,keepaspectratio]{#1}}

\sloppy

\title{Malicious Traffic Detection in CICIDS-2017 Using Machine Learning}

\author{Matheus Fraga Cardoso, Vinícius Guerra de Mello}

\address{Programa de Pós-Graduação em Computação (PPGC) -- Universidade Federal do Rio Grande do Sul (UFRGS)\\
Porto Alegre -- RS -- Brasil
}

\begin{document}
\raggedbottom

\maketitle
\pagestyle{plain}

\begin{abstract}

\end{abstract}

\begin{resumo}

\end{resumo}

\section{Introduction}

A detecção de tráfego malicioso é um problema central em cibersegurança.
Segundo o \textit{Verizon Data Breach Investigations Report} (DBIR) de 2026~\cite{verizon2026dbir}, os maiores ataques distribuídos de negação de serviço (DDoS) apresentaram crescimento de 198\% no volume de bits por segundo; a exploração de vulnerabilidades tornou-se o principal vetor inicial de invasão, representando 31\% das violações registradas e crescimento de 55\% em relação ao período anterior; e o uso de inteligência artificial por agentes maliciosos acelera a descoberta de falhas, automatiza a exploração de vulnerabilidades e potencializa etapas dos ataques cibernéticos, aumentando a complexidade dos mecanismos tradicionais de defesa.

Nesse cenário, os sistemas de detecção de intrusão (IDS) baseados em assinatura dependem de padrões conhecidos e bem definidos, mas apresentam limitações claras frente a ataques novos ou ofuscados, o que motiva abordagens estatísticas e de aprendizado de máquina (AM) capazes de generalizar para padrões não vistos.
Além disso, abordagens baseadas em AM enfrentam desafios adicionais, como a escassez de \textit{datasets} de referência publicamente disponíveis e o envelhecimento rápido desses \textit{benchmarks} à medida que as tecnologias de rede evoluem~\cite{he2023adversarial}.

Nesse contexto, este trabalho treina modelos de aprendizado de máquina para detecção binária de tráfego malicioso, a fim de obter classificadores capazes de distinguir tráfego legítimo e malicioso e de generalizar para ataques novos ou ofuscados, superando as limitações das abordagens por assinatura e mitigando os efeitos do envelhecimento de \textit{datasets}.
Para isso, utiliza o conjunto de dados CICIDS-2017~\cite{sharafaldin2018toward}, amplamente adotado como \textit{benchmark} de detecção de intrusão, com 14 tipos de ataques rotulados e mais de 2 milhões de amostras de fluxos de rede.
Este trabalho também explora a interpretabilidade dos modelos treinados e a reprodutibilidade dos experimentos propostos.

As principais contribuições deste trabalho são: (i) comparação sistemática de cinco classificadores para detecção binária no CICIDS-2017; (ii) análise de interpretabilidade do melhor modelo (SHAP, importância de atributos); (iii) pipeline reproduzível (código/Docker).

O restante deste artigo está organizado da seguinte forma.
A Seção~\ref{sec:related} revisa trabalhos relacionados à detecção de intrusão com aprendizado de máquina.
A Seção~\ref{sec:methodology} descreve a metodologia, incluindo o \textit{dataset}, o pré-processamento, o \textit{pipeline} de treinamento, as métricas de avaliação, a interpretabilidade dos modelos e a reprodutibilidade dos experimentos.
A Seção~\ref{sec:experiments} apresenta os experimentos e os resultados obtidos.
A Seção~\ref{sec:discussion} discute os achados e limitações do estudo.
Por fim, a Seção~\ref{sec:conclusion} apresenta as conclusões e perspectivas de trabalhos futuros.

\section{Related Work}
\label{sec:related}

A literatura sobre detecção de intrusão baseada em aprendizado de máquina (AM) consolidou o uso de \textit{benchmarks} públicos de tráfego rotulado para comparar classificadores e relatórios de desempenho.
Entre esses conjuntos, o CICIDS-2017~\cite{sharafaldin2018toward} tornou-se referência recorrente por reproduzir fluxos benignos e maliciosos com dezenas de atributos derivados de fluxos de rede.
Nesta seção, revisamos três trabalhos representativos: dois estudos empíricos no CICIDS-2017 e um \textit{survey} sobre AM adversarial em NIDS.

Jairu e Mailewa~\cite{jairu2022network} propõem uma abordagem supervisionada de inteligência artificial para descoberta de anomalias de rede no CICIDS-2017.
O estudo enfatiza a insuficiência de IDS baseados apenas em assinaturas e avalia múltiplos classificadores supervisionados sobre características de fluxo, com foco em identificar padrões de tráfego malicioso frente a tráfego legítimo.
Os autores discutem o papel do pré-processamento e da seleção de atributos para elevar a acurácia na detecção de ataques conhecidos no \textit{dataset}, posicionando o AM como alternativa prática à detecção puramente baseada em regras.

Maseer \textit{et al.}~\cite{maseer2021benchmarking} conduzem um \textit{benchmark} sistemático de algoritmos de AM para IDS baseados em anomalias no CICIDS2017.
O artigo compara dez técnicas supervisionadas e não supervisionadas---incluindo árvores de decisão, \textit{k}-NN, Naive Bayes, floresta aleatória, SVM e redes neurais---em dezenas de configurações de modelos, reportando acurácia, precisão, \textit{recall}, F1 e tempo de treinamento.
Os resultados indicam que modelos como \textit{k}-NN, árvore de decisão e Naive Bayes obtêm desempenho competitivo em classes de ataques específicas, mas o desbalanceamento e a diversidade de ataques no CICIDS2017 exigem critérios de avaliação cuidadosos.

Em perspectiva mais ampla, He \textit{et al.}~\cite{he2023adversarial} apresentam um \textit{survey} sobre aprendizado de máquina adversarial aplicado a sistemas de detecção de intrusão em rede (NIDS).
Os autores organizam taxonomias de NIDS baseados em aprendizado profundo, formulam ataques adversariais voltados a evasão de modelos e revisam mecanismos de defesa.
O levantamento destaca desafios estruturais da área---como escassez de \textit{datasets} públicos confiáveis, obsolescência rápida de \textit{benchmarks} e limitações de interpretabilidade em modelos de alto desempenho---que motivam avaliações reproduzíveis e análises explicáveis, e não apenas ganhos marginais de acurácia.

\textbf{Posicionamento deste trabalho.}
Os estudos de Jairu e Mailewa~\cite{jairu2022network} e de Maseer \textit{et al.}~\cite{maseer2021benchmarking} concentram-se principalmente no desempenho preditivo de classificadores no CICIDS-2017, em geral em cenários multiclasse ou fortemente orientados a métricas globais.
O \textit{survey} de He \textit{et al.}~\cite{he2023adversarial} sistematiza riscos e lacunas de pesquisa, mas não propõe um pipeline experimental único com foco em interpretabilidade.
Este trabalho diferencia-se por: (i) adotar detecção binária (tráfego legítimo \textit{versus} malicioso) com cinco classificadores clássicos; (ii) analisar interpretabilidade do melhor modelo (por exemplo, SHAP e medidas de importância de atributos); e (iii) documentar a reprodutibilidade dos experimentos (código e ambiente conteinerizado).
Assim, complementamos \textit{benchmarks} existentes no CICIDS-2017 com uma visão orientada à explicabilidade e à replicação, em linha com as recomendações de pesquisa apontadas por He \textit{et al.}~\cite{he2023adversarial}.

\section{Methodology}
\label{sec:methodology}

\subsection{Dataset}
O \textit{dataset} utilizado é o CICIDS-2017~\cite{sharafaldin2018toward}, 
produzido pelo Canadian Institute for Cybersecurity.
O conjunto compreende 14 tipos de ataques, em um cenário de classificação multiclasse 
nos rótulos originais; neste trabalho, os fluxos são agregados para detecção binária 
(tráfego legítimo \textit{versus} malicioso).

A distribuição agregada dos rótulos é severamente desbalanceada: 80{,}3\% de fluxos \texttt{BENIGN} e 19{,}7\% \texttt{ATTACK} após unificar os tipos de ataque (Figura~\ref{fig:distribuicao-binaria}).
Entre as classes de ataque individuais, a disparidade é ainda maior; categorias raras, como \texttt{Heartbleed} (7 amostras) e \texttt{Web\_Attack\_Sql\_Injection} (21 amostras), tornam inviável um treinamento multiclasse estável no escopo deste estudo (Figura~\ref{fig:distribuicao-multiclasse}).

\subsubsection{Formulação binária}
Dada a dificuldade de aprender classes com menos de 40 amostras no prazo do projeto, agrupamos todos os ataques na classe \texttt{ATTACK} e definimos a tarefa como classificação binária (\texttt{BENIGN} \textit{versus} \texttt{ATTACK}).
Reconhecemos que essa escolha reduz a granularidade por tipo de ataque; essa limitação é discutida na Seção~\ref{sec:discussion}.

\begin{figure}[!tb]
  \centering
  \includefig{../projeto/archive/eda_01_distribuicao_classes.png}
  \caption{Distribuição agregada \texttt{BENIGN} \textit{versus} \texttt{ATTACK} no CICIDS-2017.}
  \label{fig:distribuicao-binaria}
\end{figure}

\begin{figure}[!tb]
  \centering
  \includefig{../projeto/archive/eda_02_tipos_ataque.png}
  \caption{Distribuição por tipo de ataque no CICIDS-2017.}
  \label{fig:distribuicao-multiclasse}
\end{figure}

\subsection{Data and preprocessing}
\label{sec:preprocessing}

Unificamos os oito arquivos CSV do CICIDS-2017 e aplicamos um \textit{pipeline} de limpeza e redução de dimensionalidade, guiado por análise exploratória (EDA), antes do treinamento dos classificadores.

Os problemas tratados foram:
\begin{itemize}
  \item \textbf{Duplicatas:} 255.234 linhas repetidas (9{,}03\% do total), removidas com base nas colunas de fluxo (excluindo metadados de origem do arquivo).
  \item \textbf{Valores infinitos:} nas colunas \texttt{Flow Bytes/s} e \texttt{Flow Packets/s}, gerados por divisão por zero; substituídos por \texttt{NaN} e descartados na limpeza subsequente.
  \item \textbf{Variância zero:} oito atributos constantes em todo o conjunto (\texttt{*\_Avg\_Bulk\_*}, \texttt{Bwd PSH Flags}, \texttt{Bwd URG Flags}), eliminados por não contribuírem para a discriminação.
  \item \textbf{Coluna redundante:} \texttt{Fwd Header Length.1}, idêntica a \texttt{Fwd Header Length}, removida.
  \item \textbf{Alta correlação:} remoção iterativa de atributos com $|r| \geq 0{,}95$ (Pearson), resultando em 23 atributos redundantes adicionais descartados.
  \item \textbf{Outliers:} tratamento pelo método do intervalo interquartil (IQR), com limitação suave (\textit{clipping}) nos extremos das variáveis numéricas, conforme a distribuição observada na EDA (Figura~\ref{fig:outliers-eda}).
\end{itemize}

Após a agregação binária dos rótulos, a matriz final contém \textbf{47 atributos numéricos} e \textbf{2.824.951 fluxos}.

\begin{figure}[!tb]
  \centering
  \includefigtall{../projeto/archive/eda_04_boxplots_outliers.png}
  \caption{Boxplots dos atributos com maior incidência de outliers (IQR).}
  \label{fig:outliers-eda}
\end{figure}

\subsection{Training pipeline}
\label{sec:training}

O treinamento foi implementado em Python~3.11 com \textit{scikit-learn}, em cinco etapas, com \texttt{random\_state=42}, \texttt{class\_weight=balanced} (sem SMOTE) e métricas sensíveis ao desbalanceamento (Seção~\ref{sec:metrics}).

\textbf{Algoritmos avaliados.}
Comparamos cinco classificadores clássicos, escolhidos por interpretabilidade, custo computacional e uso recorrente em \textit{benchmarks} de IDS: Regressão Logística (RL), Random Forest (RF), Árvore de Decisão (AD), Hist Gradient Boosting (HGB) e Naive Bayes Gaussiano (NB).
A RL é treinada em um \texttt{Pipeline} com \texttt{StandardScaler} ajustado dentro de cada fold de validação, prevenindo \textit{data leakage} do conjunto de validação no escalonamento; árvores e ensemble baseados em árvores (RF, AD, HGB) operam diretamente sobre os atributos originais.

\textbf{Divisão dos dados.}
Aplicamos \textit{holdout} estratificado 80/20 sobre a matriz final de fluxos, preservando a proporção \texttt{BENIGN}/\texttt{ATTACK} em treino e teste.
O conjunto de teste permaneceu isolado de todas as etapas subsequentes---validação cruzada, testes estatísticos e otimização de hiperparâmetros---e foi utilizado apenas na avaliação final reportada.

\textbf{Experimentos.}
\begin{enumerate}
  \item \textbf{\textit{Spot-checking}:} treino dos cinco modelos com hiperparâmetros padrão da biblioteca (por exemplo, RL com \texttt{max\_iter=1000}, RF com \texttt{n\_estimators=100}) sobre o conjunto de treino e avaliação no conjunto de teste, para triagem inicial.
  \item \textbf{Validação cruzada:} os mesmos modelos foram reavaliados por \textit{Repeated Stratified K-Fold} com $k=10$ \textit{folds} e três repetições (30 avaliações por modelo), exclusivamente no conjunto de treino, reportando média e desvio padrão de F1, precisão, \textit{recall} e ROC-AUC (classe positiva \texttt{ATTACK}).
  \item \textbf{Testes estatísticos:} teste de Friedman como procedimento \textit{omnibus} sobre os 30 escores por modelo, seguido de Wilcoxon pareado com correção de Bonferroni entre pares de classificadores.
  \item \textbf{Otimização de hiperparâmetros:} para os dois modelos mais promissores nas etapas anteriores (RF e HGB), \textit{RandomizedSearchCV} com 20 combinações amostradas e \textit{StratifiedKFold} interno com cinco \textit{folds}; por viabilidade computacional, a busca usou uma amostra estratificada de 20\% do treino, e o melhor modelo foi retreinado no treino completo.
  \item \textbf{Avaliação final:} modelo otimizado avaliado no \textit{holdout} de teste reservado.
\end{enumerate}

\subsection{Metrics and Evaluation}
\label{sec:metrics}

Avaliamos os classificadores com quatro métricas calculadas no \textit{scikit-learn}, sempre com classe positiva \texttt{Attack} (tráfego malicioso): precisão (fração de alertas corretos entre os previstos como ataque), \textit{recall} (fração de ataques reais detectados), F1 (média harmônica entre precisão e \textit{recall}) e ROC-AUC (capacidade de separar as classes ao variar o limiar de decisão).
Também reportamos acurácia (\textit{accuracy}) como referência, mas não a usamos para ranquear modelos.

As classes previstas vêm de \texttt{predict()}, que no \textit{scikit-learn} aplica limiar 0{,}5 sobre $P(\texttt{Attack})$ quando o modelo fornece probabilidades.
No \textit{spot-checking} e na avaliação final, precisão, recall e F1 usam \texttt{pos\_label=\"Attack\"} com rótulos categóricos \texttt{Benign}/\texttt{Attack}; o ROC-AUC usa \texttt{predict\_proba} da classe \texttt{Attack} (métrica independente de limiar).
Na validação cruzada, codificamos $y$ como binário ($1=\texttt{Attack}$, $0=\texttt{Benign}$) e aplicamos os \textit{scorers} \texttt{f1}, \texttt{precision}, \texttt{recall} e \texttt{roc\_auc}: os três primeiros medem a classe positiva codificada como 1 (comportamento binário focado em \texttt{Attack}, sem \texttt{average='macro'} ou \texttt{'weighted'}).

Adotamos o F1 da classe \texttt{Attack} como critério principal de comparação e seleção dos modelos.
Com cerca de 80\% de fluxos \texttt{Benign} e 20\% \texttt{Attack}, um classificador que prevê sempre tráfego legítimo obtém acurácia elevada, porém F1 nulo para ataques---inadequado para detecção de intrusão, em que falhar em detectar \texttt{Attack} é o erro mais grave.
O F1 equilibra precisão e \textit{recall}: alta precisão reduz falsos alarmes; alto \textit{recall} aumenta a cobertura de ataques reais.

No \textit{spot-checking}, calculamos as cinco métricas no conjunto de teste reservado, geramos matriz de confusão e \textit{report} de classificação por modelo, e ordenamos os resultados pelo F1 de \texttt{Attack}.
Na validação cruzada (\textit{Repeated Stratified K-Fold}, 30 avaliações por algoritmo), registramos média e desvio padrão de F1, precisão, \textit{recall} e ROC-AUC apenas no conjunto de treino; esses escores alimentaram os testes de Friedman e Wilcoxon.
Após a otimização de hiperparâmetros (RF e HGB), a avaliação final no \textit{holdout} de teste repetiu o mesmo conjunto de métricas para estimar o desempenho de generalização.

\subsection{Interpretability}
\label{sec:interpretability}

\FloatBarrier

Além do desempenho preditivo, analisamos \textit{por que} o classificador separa tráfego legítimo e malicioso---requisito central em IDS, em que operadores precisam auditar alertas e validar se o modelo aprendeu sinais de rede plausíveis.
Tomamos como referência o \textbf{Random Forest otimizado} (\texttt{rf\_opt}), retreinado no conjunto de treino completo após \textit{RandomizedSearchCV}: no \textit{holdout}, obteve o maior F1 da classe \texttt{Attack} entre os modelos ajustados (RF e HGB), além de ser naturalmente interpretável por árvores.
Todas as análises abaixo usam os mesmos \textbf{47 descritores numéricos} produzidos na Seção~\ref{sec:preprocessing} (espaço de atributos de treino, teste e SHAP).

\subsubsection{Importância global (Gini)}

Extraímos \texttt{feature\_importances\_} do \texttt{rf\_opt} já ajustado em $X_{\mathrm{train}}$: em cada árvore, agrega-se a redução média de impureza de Gini por atributo, gerando um ranqueamento global para triagem operacional.
Esse indicador é rápido, mas pode inflar atributos com alta cardinalidade ou fortemente correlacionados; por isso o complementamos com explicações locais (SHAP).
A Figura~\ref{fig:interp-gini} mostra o top-20; lideram \texttt{Destination Port} e descritores \textit{backward}.

\begin{figure}[!tb]
  \centering
  \includefigtall{../projeto/archive/interp_01_feature_importance_gini.png}
  \caption{Top-20 atributos por importância de Gini no \texttt{rf\_opt} (treino completo, 47 descritores).}
  \label{fig:interp-gini}
\end{figure}

\subsubsection{Explicações locais (SHAP)}

O protocolo SHAP é fixo e reproduzível: o modelo é o \texttt{rf\_opt} retreinado em $X_{\mathrm{train}}$ completo (sem novo ajuste de hiperparâmetros); as contribuições são calculadas com \texttt{TreeExplainer}~\cite{lundberg2017unified} sobre uma amostra aleatória de \textbf{5.000} fluxos de $X_{\mathrm{test}}$ (\texttt{random\_state=42}), para a classe \texttt{Attack} (valores SHAP positivos favorecem \texttt{Attack}; negativos, \texttt{Benign}).
Trata-se de explicação \textbf{pós-hoc}: não há retreino do classificador nem seleção ou ajuste de atributos com base no conjunto de teste---o SHAP apenas quantifica, para o modelo já fixo, quanto cada um dos 47 descritores altera a predição em cada fluxo amostrado.

\begin{figure}[!tb]
  \centering
  \includefigtall{../projeto/archive/interp_02_shap_beeswarm.png}
  \caption{Resumo SHAP (\textit{beeswarm}), classe \texttt{Attack} ($n=5.000$, teste).}
  \label{fig:interp-shap-beeswarm}
\end{figure}

A Figura~\ref{fig:interp-shap-beeswarm} resume o efeito por atributo: \texttt{Destination Port} e estatísticas \textit{backward} concentram contribuições positivas para \texttt{Attack}, alinhadas ao ranqueamento de Gini.
Para as cinco variáveis de maior $|SHAP|$ médio, a Figura~\ref{fig:interp-shap-dep} traz \textit{dependence plots} (efeito marginal e não linearidades).

\begin{figure}[!tb]
  \centering
  \includefigtall{../projeto/archive/interp_03_shap_dependence.png}
  \caption{\textit{Dependence plots} SHAP: top-5 por $|SHAP|$ médio (classe \texttt{Attack}).}
  \label{fig:interp-shap-dep}
\end{figure}

No top-20 de Gini e no ranqueamento por $|SHAP|$ médio, o núcleo de atributos coincide (\texttt{Destination Port}, descritores \textit{backward} e tamanhos de pacote), indicando concordância qualitativa entre importância global e local.

\paragraph{Limitações.}
A importância de Gini não desagrega efeitos entre atributos correlacionados; o SHAP foi avaliado em 5.000 fluxos de teste, não no \textit{holdout} inteiro; \texttt{TreeExplainer} aplica-se a modelos baseados em árvore (aqui, RF) e não se estende, sem adaptação, a outros classificadores; as conclusões interpretativas limitam-se ao CICIDS-2017 e à partição treino/teste deste estudo.

\FloatBarrier

\subsection{Reproducibility}
\label{sec:reproducibility}

O pipeline está em \texttt{projeto/códigos/notebook\_completo.ipynb} no repositório \url{https://github.com/vinigm/projeto-ddos} (commit \texttt{2b0b79d}), com caminhos relativos a \texttt{../archive/} (figuras, \texttt{raw\_limpo.parquet}) e \texttt{../../artifacts/} (cache).
O \textit{parquet} pré-processado (47 atributos; Seção~\ref{sec:preprocessing}) é baixado da release v1.0 se não existir localmente.
Todas as etapas estocásticas usam \texttt{random\_state=42}; dependências fixadas em \texttt{requirements.txt} (Python~$\geq$3.10; \texttt{scikit-learn==1.8.0}, \texttt{shap==0.51.0}, entre outras).
Com \texttt{RECOMPUTE=False}, modelos, escores de CV, hiperparâmetros e valores SHAP são lidos de \texttt{artifacts/} quando disponíveis; \texttt{RECOMPUTE=True} refaz o treino completo.
Para replicar resultados e figuras, use \textit{Run All Below} a partir da seção ``Pós EDA'' do notebook (com \texttt{RECOMPUTE=False} quando o cache existir); o guia \texttt{DOCKER.md} descreve a execução com \texttt{docker compose up --build}~\cite{boettiger2015docker}.
Artefatos volumosos podem não estar versionados; a primeira replicação pode exigir gerá-los localmente.

\FloatBarrier
\clearpage

\section{Experiments and Results}
\label{sec:experiments}

Avaliamos cinco classificadores (RL, RF, AD, HGB e NB) nas cinco etapas da Seção~\ref{sec:training}.
No \textit{spot-checking} com hiperparâmetros padrão, ensembles baseados em árvores (RF, AD e HGB) superaram RL e NB no F1 da classe \texttt{Attack} no \textit{holdout} reservado; o Naive Bayes Gaussiano ficou aquém, coerente com atributos contínuos e desbalanceamento.

Na validação cruzada (\textit{Repeated Stratified K-Fold}, 30 avaliações por modelo), o RF obteve o maior F1 médio para \texttt{Attack} ($0{,}9973 \pm 0{,}0001$), seguido de AD ($0{,}9971 \pm 0{,}0001$), HGB ($0{,}9969 \pm 0{,}0004$), RL ($0{,}8640 \pm 0{,}0012$) e NB ($0{,}5753 \pm 0{,}0024$).
O ROC-AUC médio acompanhou a mesma ordem: RF ($0{,}9998$), HGB ($0{,}9999$), AD ($0{,}9984$), RL ($0{,}9818$) e NB ($0{,}7265$).
Após \textit{RandomizedSearchCV}, o RF otimizado (\texttt{rf\_opt}) manteve desempenho elevado no conjunto de teste e foi selecionado para a análise de interpretabilidade (Seção~\ref{sec:interpretability}).

\section{Discussion}
\label{sec:discussion}

Os resultados indicam que modelos baseados em árvores capturam bem a separação binária no CICIDS-2017 após o pré-processamento descrito, enquanto o NB Gaussiano permanece inadequado como \textit{baseline} probabilístico neste cenário.
O \texttt{rf\_opt} equilibra alto F1 de \texttt{Attack} e interpretabilidade operacional: as importâncias de Gini e os valores SHAP convergem para \texttt{Destination Port} e estatísticas de fluxo \textit{backward}, sinais plausíveis em IDS.

A formulação binária simplifica o problema e mascara ataques raros (por exemplo, \texttt{Heartbleed} e \texttt{Web\_Attack\_Sql\_Injection}); trabalhos futuros podem explorar cenários multiclasse com técnicas de balanceamento ou detecção hierárquica.
Também permanece em aberto a avaliação de ataques específicos pouco frequentes e a robustez frente a evasão adversarial, tema central no \textit{survey} de He \textit{et al.}~\cite{he2023adversarial}.
Por fim, investigar LLMs para auxiliar geração de tráfego sintético adversarial---ou para priorizar atributos em análise---pode complementar, e não substituir, descritores de fluxo tabulares como os do CICIDS-2017.

\section{Conclusion}
\label{sec:conclusion}

Este trabalho comparou cinco classificadores para detecção binária de tráfego malicioso no CICIDS-2017, com pipeline reproduzível e análise de interpretabilidade do RF otimizado.
O RF obteve o melhor F1 médio na validação cruzada e foi corroborado por explicações globais (Gini) e locais (SHAP) alinhadas a descritores de rede conhecidos.
Concluímos que ensembles de árvores são candidatos fortes a IDS baseados em AM neste \textit{benchmark}, mas a detecção prática ainda exige endereçar desbalanceamento extremo por tipo de ataque, custo computacional em escala e validação em tráfego fora do conjunto de treino.

\bibliographystyle{sbc}
\bibliography{referencias}
% Para os trabalhos que envolverem dados publicamente disponíveis, os grupos
% devem incluir o link para download no seu artigo.

\end{document}
